% !TEX root = ../print.tex
% Draft \S14 --- Implementation: the hemigroup is the algorithm
%
% Source: draft/spatial-hemigroup-scale-space.md, \S14 (lines 823-876).
%
% SCOPE. As the causal article does for the same two chapters, the nodes here and in
% Chapter 15 are blueprint statements and not formalisation targets: none carries a \lean
% tag and none is scheduled in a Lean phase. The blueprint is nevertheless complete with
% respect to the article's mathematics.
%
% CHANGED (prop:rational-increments clause 2): the draft writes the interpolated profile as
% sum_{m >= 0} k_R(q^m x). It is sum_{m >= 1}: with F = sum_{m >= 1} psi(q^{-m} .) the
% dilate psi(c .) has profile k_R(u/c), so c = q^{-m} contributes k_R(q^m u) for m >= 1, and
% the m = 0 term would make the q-step increment psi(q .) rather than psi. The node states
% the corrected form; the draft should follow.
%
% NARROWED (prop:rational-increments): the draft's clause (3), that an admissible family all
% of whose increments are rational of bounded order has the form of clause (1), is a sketch
% with three unresolved steps of its own (in what sense the difference quotients converge,
% whether bounded order is removable, and the exclusion of oscillating terms). It is not
% stated as a node; it is recorded as an open statement in rem:rational-converse.
%
% CHANGED (prop:thorin-machine): the draft's parallel generator is written
% "2a d_x^2 + sum_i w_i (L_{t/theta_i} - I) . (2/t)". Both constants are wrong against
% prop:corner-generators: the Gaussian term is 2at d_x^2 and the sum carries 1/t, not 2/t.
% The node states the corrected form.
%
% (Checked at the draft sync, 2026-09-11: the draft already carried both corrected forms above --
% the profile sum from m >= 1 and the generator with 2at and 1/t -- so those two CHANGED notes
% record a past state of the draft. The same sync narrowed prop:thorin-machine's rounding to even
% weights, R82, where the draft had been right and the node wrong.)
%
% The draft's dependence on the pyramid module ([VI, ...]) is removed: the q-difference
% rigidity used in clause (2) is proved inline in three lines, so nothing here waits on an
% unreleased module.

\chapter{Implementation: the hemigroup is the algorithm}
\label{ch:implementation}

The structural fact that organizes everything: the scale direction never needs to be
discretized approximately. What is new on the line against the causal theory is that the
causal first-order section becomes a forward--backward pair --- the recursive filter the
image-processing literature has used for the Gaussian for decades --- and that the theory says
which of those filters are scale spaces and which are not.

% draft: Proposition 14.1
% twin: hcs prop:knot-exactness --- verbatim.
\begin{proposition}[exactness at the knots]\label{prop:knot-exactness}
\uses{lem:convolution-representation,def:cascade-family,lem:additivity}
\notes{exact-cascade-implementation}
Let $0 = t_0 < t_1 < \dots < t_M$ be any finite set of scale knots and define the discrete
cascade $u_0 := f$, $u_{k+1} := \mu_{t_k,t_{k+1}} * u_k$. Then $u_k = \Phi_{0,t_k}f$ exactly for
every $k$, and refining the knot set changes nothing at the old knots. On the transform side the
increment transforms telescope,
$\prod_{k<M} e^{-F(t_{k+1}\omega) + F(t_k\omega)} = e^{-F(t_M\omega)}$.
\statusT\quad\emph{(One line of hemigroup algebra; the content is that the cascade is not an
approximation scheme for the family but the family itself, sampled. Since the family exists at
every scale, a ladder of knots can be shifted at will, an arbitrary dilation being handled by
instantiating the family at the shifted knots.)}
\end{proposition}

\begin{proof}
Induction on (A6): $\mu_{0,t_{k+1}} = \mu_{t_k,t_{k+1}} * \mu_{0,t_k}$, by
\ref{lem:additivity}, and $u_0 = \Phi_{0,t_0}f = f$.
\end{proof}

% draft: Proposition 14.2
% twin: hcs prop:gamma-cascade --- the causal section is one-sided; here it is a
% forward--backward pair, and the increment kernel is the four-term mixture (14.2).
% CHANGED (2026-09-19, article mirror of module B, R174): four register repairs in the statement,
% from the register pass's items 4, 5 and 9. None changes what is claimed except the last, which
% narrows.
% (a) The term LAG SECTION is introduced here, at the first place the sections are constructed;
% chapter 14 and prop:rational-closure used it with no gloss, and the paper introduces it in its
% own prose ("called a lag section below").
% (b) "free of input differentiation" is replaced by what it says: the zero of the transfer
% function costs no differentiation of the input, the output being a convex combination of the
% input and the state. The same sentence is added to the proof, where the convex combination is
% computed.
% (c) Expo_+- IS DEFINED, and \Lap takes the subscript it takes in chapters 11, 13 and 15
% (\Lap(t) here was the article's only argument form). eq:four-term reads \Lap_{t_{k+1}}.
% (d) THE COST IS NARROWED to a sampled realization: the section is an ODE in x, so "multiply-adds
% per sample" presupposes a grid, which the statement did not state. The sentence now says so and
% says that the sampling of the signal is not treated here, as the paper's opening paragraph does.
% CHANGED 2026-09-19 (R183, second external review round, referee C finding 11): TWO CLARIFICATIONS
% in the statement, neither of them a new claim. (a) WHICH SOLUTION of the state equation is the
% realization: its solutions differ by a multiple of e^{-x/t_{k+1}}, and the transfer identity of
% the proof holds for the L^1 one (w(-infty) = 0 forward, w(+infty) = 0 backward). On a finite or
% periodic record this is the initialization problem, which the article treats only inside the
% numerical example. (b) "IN TWO PASSES" IS PER KNOT: each record u_{k+1} needs its own forward and
% backward pass, so a ladder of M knots is 2M passes, which the O(gamma M) count already reflects.
% The per-sample multiply-add constant of rem:costs is corrected with it (2 gamma -> 4 gamma).
\begin{proposition}[the \Matern{} cascade: forward--backward first-order sections]
\label{prop:matern-cascade}
\uses{prop:knot-exactness,prop:matern-exponent,lem:nonvanishing}
\notes{exact-cascade-implementation}
For the \Matern{} family with integer $\gamma$, so that $e^{-F(t\omega)} =
(1+t^2\omega^2)^{-\gamma}$, the increment between adjacent knots has transform
\begin{equation}\label{eq:matern-increment}
  \hat\mu_{t_k,t_{k+1}}(\omega)
  = \Bigl(\frac{1 + t_k^2\omega^2}{1 + t_{k+1}^2\omega^2}\Bigr)^{\gamma}
  = \Bigl(\frac{1 + it_k\omega}{1 + it_{k+1}\omega}\cdot
          \frac{1 - it_k\omega}{1 - it_{k+1}\omega}\Bigr)^{\gamma},
\end{equation}
a cascade of $\gamma$ identical causal first-order pole--zero sections in $x$, called \emph{lag
sections} below, and $\gamma$ identical anticausal ones, each with the one-state realization
\[
  w' = \frac{y - w}{t_{k+1}} \quad(\text{forward}),
  \qquad
  v = \frac{t_k}{t_{k+1}}\,y + \Bigl(1 - \frac{t_k}{t_{k+1}}\Bigr) w,
\]
and its mirror image run backward. Which solution of the state equation is meant is part of the
realization: its solutions differ by a multiple of $e^{-x/t_{k+1}}$, and the one intended is the
solution in $\Lone$, equivalently the one with $w(-\infty) = 0$ for the forward pass and
$w(+\infty) = 0$ for the backward one. The zero of
the transfer function costs no differentiation of
the input: the output is a convex combination of the input and the state. The section's kernel is
the mixture
$\tfrac{t_k}{t_{k+1}}\delta_0 + (1 - \tfrac{t_k}{t_{k+1}})\,\Expo_{+}(t_{k+1})$ of no
displacement and a one-sided exponential displacement, $\Expo_{\pm}(t)$ denoting the one-sided
exponential kernels $t^{-1}e^{-|x|/t}$ carried by $\pm(0,\infty)$; so the two-sided increment for
$\gamma = 1$ is the four-term mixture
\begin{equation}\label{eq:four-term}
  r^2\,\delta_0 + r(1-r)\bigl(\Expo_+ + \Expo_-\bigr) + (1-r)^2\,\Lap_{t_{k+1}},
  \qquad r := \frac{t_k}{t_{k+1}} \in (0,1),
\end{equation}
where $\Expo_\pm$ are taken at the range $t_{k+1}$ and $\Lap_t$ is the Laplace kernel of range
$t$; positivity and unit mass are manifest at the
realization level. Sampled on a grid, a full scale space over $M$ knots costs $2\gamma M$
first-order states and $O(\gamma M)$ multiply--adds per sample, in a forward and a backward pass
\emph{at each knot}. Sampling the signal
is a separate approximation and is not treated here.
\statusT
\end{proposition}

\begin{proof}
The increment transform is the ratio of the endpoint transforms, the denominators being
nonvanishing (\ref{lem:nonvanishing}), which is \ref{prop:matern-exponent}(2); the
factorization is $1 + t^2\omega^2 = (1+it\omega)(1-it\omega)$. The transfer identity of the
section is
\[
  \frac{t_k}{t_{k+1}} + \Bigl(1 - \frac{t_k}{t_{k+1}}\Bigr)\frac{1}{1+it_{k+1}\omega}
  = \frac{1 + it_k\omega}{1 + it_{k+1}\omega},
\]
and the convex-combination clause is $0 < t_k/t_{k+1} < 1$. In particular the output is a convex
combination of the input and the state, so the section differentiates nothing. For
\ref{eq:four-term}, convolve the
causal and anticausal mixtures and use $\Expo_+(t) * \Expo_-(t) = \Lap_t$, whose transform identity
is $(1+it\omega)^{-1}(1-it\omega)^{-1} = (1+t^2\omega^2)^{-1}$.
\end{proof}

% CHANGED (2026-09-19, article mirror of module B, R174): THE CHAPTER'S CHECK IS DISCHARGED (the
% register pass's item 17). The two standard recursive Gaussians are named, with the pole structure
% the librarian read from held copies on 2026-09-18: Deriche, INRIA RR-1893, section 5.1, eq. (38),
% two complex-conjugate pole pairs; Young and van Vliet 1995, p. 141, Table 1 with (6a)-(6b), one
% real pole and one conjugate pair. Both therefore fall in the second half of the remark, which is
% about complex poles, and the paper's example applies the test to both. Two further repairs in the
% same pass: the infinite-divisibility criterion for a rational transfer, which was asserted with no
% pointer, is now attributed to the uniqueness of the \Levy{} pair, which is what makes the signed
% density the only candidate; and the sentence about a dominant real pole carrying a complex pair,
% which read as if real-pole members were the only admissible rational filters, points at
% rem:rational-converse, where such a family is exhibited (R170, the same day).
% NEW NODE 2026-09-19 (R198; the author's decision D4, a widening authorized at the second
% external review round of module B). The criterion deciding which rational transfer functions
% are admissible kernels was a remark -- rem:recursive-gaussians, which stays below as a short
% remark pointing here, so that existing references survive -- and three things were wrong with
% it as a remark (referee C findings 4 and 5, referee E section 2.7). Its opening sentence was
% circular ("R is the kernel of an admissible family exactly when R = e^{-F} for an admissible
% F"). Its frame, R = |P/Q|^2, does not cover a design that ADDS a causal and an anticausal half,
% which is Deriche's; the right frame is an even rational R with R(0) = 1. And the uniqueness step
% it invoked is stated for POSITIVE \Levy{} pairs, while what is compared here is a signed density
% -- the recombination of lem:cin-delay-equation's proof is what closes the gap, and it is written
% out below. The identity of prop:matern-exponent(1) at a complex rate, asserted in
% rem:rational-converse and used there and here, is proved here, at its first use.
% [T], prose-only: no Lean is scheduled for this chapter and none is affected.
\begin{proposition}[which rational transfer functions are admissible kernels]
\label{prop:rational-criterion}
\uses{thm:main-characterization,lem:selfdecomposable-exponents,lem:profile-integrability,prop:fourier-toolbox,prop:matern-exponent,lem:nonvanishing}
\notes{exact-cascade-implementation}
Let $R$ be an even rational function with $R(0) = 1$, written in lowest terms. Write
$\pm i\theta_1,\dots,\pm i\theta_M$ for its poles and $\pm i\zeta_1,\dots,\pm i\zeta_L$ for its
zeros, listed with multiplicity and normalised by $\Rea\theta_m > 0$ and $\Rea\zeta_l > 0$; both
lists are closed under conjugation, so the \emph{signed profile}
\begin{equation}\label{eq:rational-profile}
  k(x) \ :=\ 2\sum_{m \le M} e^{-\theta_m x} \ -\ 2\sum_{l \le L} e^{-\zeta_l x},
  \qquad x > 0,
\end{equation}
is real-valued. Then $R$ is the transform of the kernel at canonical scale $1$ of an admissible
family --- equivalently, $-\log R$ is an admissible exponent --- if and only if
\begin{enumerate}
\item $R$ has neither a pole nor a zero on the real axis, so that $R > 0$ there;
\item $k \ge 0$ on $(0,\infty)$; and
\item $k$ is nonincreasing on $(0,\infty)$.
\end{enumerate}
The family is then the one with Gaussian coefficient $0$ and folded profile $k$. Condition (i)
alone excludes every $R$ with a real zero, since the transform of an infinitely divisible law has
no zero at all (\ref{lem:nonvanishing}).
\statusT\quad\emph{(The test a practitioner applies to a recursive filter, and the reason this
chapter can say which published designs are scale spaces and which are not. The three conditions
are read off the pole--zero diagram: (i) is a glance at the real axis, and (ii) and (iii) are two
inequalities for one explicit exponential sum. A Gaussian factor is not covered and needs none: a
rational $R$ forces the Gaussian coefficient to vanish, which is the growth argument in the proof,
and a design with a Gaussian factor is that factor times a rational one. \textbf{Not claimed}: any
statement about a \emph{sampled} filter --- the criterion is for the transfer function on the
line, and the discrete classification is a different theory --- and no bound on how far a failing
filter is from the class, which is a numerical question. The formula \ref{eq:rational-profile} is
the folded profile, so its factor $2$ is this article's convention; at a real rate $\theta$ it is
the profile $2e^{-\theta x}$ of \ref{prop:matern-exponent}(1), whose exponent is
$\log(1 + \omega^2/\theta^2)$.)}
\end{proposition}

\begin{proof}
\emph{The factorisation.} $R$ is even and rational, so $R(\omega) = N(\omega^2)/D(\omega^2)$ for
polynomials $N$ and $D$ with no common factor, real because $R$ is real on the real axis. Its
poles and zeros are therefore closed under $\omega \mapsto -\omega$ and under conjugation, and
under (i) none of them is real; choosing from each pair $\pm\omega_0$ the one in the upper
half-plane and writing it as $i\theta$ with $\Rea\theta > 0$ gives the two lists of the statement,
each closed under conjugation. Then $D(\omega^2) = c_D\prod_m(\omega^2 + \theta_m^2)$ and
$N(\omega^2) = c_N\prod_l(\omega^2 + \zeta_l^2)$, and $R(0) = 1$ fixes the constant:
\begin{equation}\label{eq:rational-factored}
  R(\omega) \ =\ \frac{\prod_{l \le L}\bigl(1 + \omega^2/\zeta_l^2\bigr)}
                      {\prod_{m \le M}\bigl(1 + \omega^2/\theta_m^2\bigr)} .
\end{equation}

\emph{The exponential-profile identity at a complex rate.} For $\Rea\theta > 0$ and $\omega$ real,
\begin{equation}\label{eq:complex-frullani}
  \int_0^\infty \bigl(1 - \cos\omega x\bigr)\,e^{-\theta x}\,\frac{dx}{x}
   \ =\ \tfrac12\bigl[\Log(\theta - i\omega) + \Log(\theta + i\omega)\bigr] - \Log\theta ,
\end{equation}
with $\Log$ the principal logarithm, which is what \ref{prop:matern-exponent}(1) says at a real
rate. Indeed $1 - \cos\omega x = 1 - \tfrac12(e^{i\omega x} + e^{-i\omega x})$, so the integral is
half the sum of $\int_0^\infty(e^{-\theta x} - e^{-(\theta \mp i\omega)x})\,dx/x$, and for two
points $\alpha,\beta$ of the open right half-plane
\[
  \int_0^\infty \bigl(e^{-\alpha x} - e^{-\beta x}\bigr)\frac{dx}{x}
   = \int_0^\infty\!\!\int_{[\alpha,\beta]} e^{-sx}\,ds\,dx
   = \int_{[\alpha,\beta]}\frac{ds}{s} = \Log\beta - \Log\alpha ,
\]
the segment $[\alpha,\beta]$ lying in that half-plane by convexity: the exchange is Fubini, the
modulus of the inner integrand being at most $e^{-cx}$ with $c > 0$ the least real part on the
segment, which is integrable against $dx$ on $[1,\infty)$ while on $(0,1)$ the inner integral is
bounded by $|\beta - \alpha|$; and $1/s$ is analytic on a neighbourhood of the segment, where
$\Log$ is a primitive.

\emph{The exponent of the signed profile.} $k$ of \ref{eq:rational-profile} is continuous on
$(0,\infty)$, bounded there by $2(M+L)$, and $|k(x)| \le 2(M+L)e^{-cx}$ with
$c := \min_m\Rea\theta_m \wedge \min_l\Rea\zeta_l > 0$; so $\int_0^\infty(1 - \cos\omega x)
|k(x)|\,dx/x$ is finite, the integrand being $O(x)$ at the origin and exponentially small at
infinity. Summing \ref{eq:complex-frullani} over the two lists,
\[
  \Psi(\omega) \ :=\ \int_0^\infty \bigl(1 - \cos\omega x\bigr)\,k(x)\,\frac{dx}{x}
   \ =\ \sum_{m}\Log\Bigl(1 + \frac{\omega^2}{\theta_m^2}\Bigr)
        - \sum_{l}\Log\Bigl(1 + \frac{\omega^2}{\zeta_l^2}\Bigr) ,
\]
where each summand is $\Log(\theta-i\omega) + \Log(\theta+i\omega) - 2\Log\theta$ and no branch is
chosen for a product: the value is real, since conjugate rates contribute conjugate summands and a
real rate a real one. Exponentiating kills the branches, $e^{\Log z} = z$ for every $z \ne 0$, and
\ref{eq:rational-factored} gives $e^{-\Psi(\omega)} = R(\omega)$. So $\Psi = -\log R$ with the
real logarithm of a positive number, at every real $\omega$.

\emph{Sufficiency.} Assume (i), (ii) and (iii). Then $k$ is nonnegative, nonincreasing, bounded on
$(0,1]$ and exponentially small at infinity, so $\int_0^1 x\,k(x)\,dx$ and
$\int_1^\infty k(x)\,x^{-1}dx$ are finite: by \ref{lem:profile-integrability} the pair $(0,k)$ is
of the form \ref{eq:sd-profile}, and by \ref{thm:main-characterization} with clause (3) of
\ref{lem:selfdecomposable-exponents} it is the pair of an admissible family, whose exponent is
$\Psi = -\log R$. The transform of its kernel at canonical scale $1$ is $e^{-\Psi} = R$.

\emph{Necessity.} Assume $R = e^{-F}$ with $F$ admissible, of data $(a,k_F)$. Then $R > 0$ and
$R \le 1$ on the real axis, so $R$ has there neither a zero nor a pole, which is (i). Next
$a = 0$: by \ref{eq:rational-factored} $|R(\omega)|$ is bounded above and below by multiples of
$|\omega|^{2(L-M)}$ for large $|\omega|$, so $F = -\log R$ grows at most like a multiple of
$\log|\omega|$, while $F(\omega) \ge a\omega^2$ for every $\omega$. So $F = \Psi$ by the previous
paragraph, that is
\[
  \int_0^\infty (1 - \cos\omega x)\,k_F(x)\,\frac{dx}{x}
   = \int_0^\infty (1 - \cos\omega x)\,k(x)\,\frac{dx}{x}, \qquad \omega \in \RR ,
\]
an identity between an exponent with a positive \Levy{} measure and one with a signed density.
Uniqueness of the \Levy{} pair (\ref{prop:fourier-toolbox}(3)) is stated for positive measures, so
it is applied after a recombination, as in the proof of \ref{lem:cin-delay-equation}: writing
$k = k^{+} - k^{-}$ for the positive and negative parts, both $(k_F + k^{-})(x)\,dx/x$ and
$k^{+}(x)\,dx/x$ are positive measures on $(0,\infty)$ satisfying the \Levy{} condition --- each
of $k^{\pm}$ is dominated by $|k|$, which does --- and the displayed identity says that the two
pairs $(0, (k_F+k^{-})\,dx/x)$ and $(0, k^{+}\,dx/x)$ have the same exponent. Hence they are
equal, and $k_F = k$ almost everywhere on $(0,\infty)$. Since $k$ is continuous and $k_F$ is
nonnegative and nonincreasing, the identity holds everywhere and $k$ inherits both properties:
the agreement set is conull, hence meets every interval, and approaching a given $x$ from each
side through it gives $k(x) \ge k(y)$ for $x < y$ and $k \ge 0$. That is (ii) and (iii).

\emph{The last sentence.} If $R$ vanishes at a real $\omega_0$ then $R$ is the transform of no
infinitely divisible law, the transform of such a law having no zero (\ref{lem:nonvanishing}); in
particular it is the transform of no admissible kernel.
\end{proof}

% CHANGED 2026-09-19 (R198, the author's decision D4): the remark KEEPS ITS LABEL and becomes
% short, its criterion being prop:rational-criterion above. What stays here is what is a remark:
% the reading of the criterion at the real-pole designs, the two published recursive Gaussians
% with the pole structure the librarian read from the held copies on 2026-09-18, and the
% recommendation. The circular opening sentence and the |P/Q|^2 frame are gone with it, and the
% two designs are now decided by the numbered conditions, Deriche's by (i) and Young and
% van Vliet's by (ii).
\begin{remark}[which recursive Gaussians are scale spaces]
\label{rem:recursive-gaussians}
\uses{prop:rational-criterion}
Read at the designs of the image-processing literature, \ref{prop:rational-criterion} says the
following. If the transfer has only real poles and no zeros, then
$R = \prod_i(1+t_i^2\omega^2)^{-1}$, whose signed profile $2\sum_i e^{-x/t_i}$ is positive and
decreasing, so the three conditions hold and the filter is the kernel of the symmetric Thorin
member with atoms at $1/t_i$ (\ref{prop:thorin-subclass}), the multi-range \Matern{}; its dilates
form a scale space whose increments are again forward--backward lag sections. The classical
all-pole approximations of the Gaussian with real poles --- in particular the cascade of $\gamma$
identical first-order sections at $t = \sigma/\sqrt{2\gamma}$, whose transfer is
$(1 + \sigma^2\omega^2/2\gamma)^{-\gamma} \to e^{-\sigma^2\omega^2/2}$ --- are therefore not
approximations of a scale space but scale spaces, \Matern{} members converging to the Gaussian
along $\gamma$ at matched variance, the variance being $2\gamma t^2$ by
\ref{prop:matern-exponent}(3). Neither complex poles nor real zeros are excluded in advance, and
neither is harmless. A complex pair contributes a term $e^{-cx}\cos(dx)$ to the signed profile,
which changes sign by itself, so condition (ii) has to be checked against the other terms; a real
pole of sufficient weight can carry such a pair, and \ref{rem:rational-converse} exhibits an
admissible family that does, of exponent $3\log(1+\omega^2) + \log(1+\omega^4/4)$. A real zero
\emph{of the transfer} fails (i) outright, while a zero off the real axis is a term subtracted in
\ref{eq:rational-profile} and may leave (ii) and (iii) standing or not: the second example of
\ref{rem:rational-converse}, of exponent $2\log(1+\omega^2) - \log(1+\omega^2/4)$, has all its
poles and zeros on the imaginary axis, is admissible, and is neither a Thorin member nor realized
by lag sections. The two standard recursive Gaussians of the image-processing literature are
optimized for accuracy and are decided by different conditions. The design of
\sscite{deriche1993recursively} \emph{adds} a causal and an anticausal half (\S5.1, (38), two
complex-conjugate pole pairs), so its transfer is not a modulus squared at all, and at the
published coefficients that sum has zeros on the real axis: it fails (i). The design of
\sscite{young1995recursive} is a forward--backward cascade with one real pole and one conjugate
pair (p.~141, Table~1 with (6a)--(6b)); it passes (i), and at the published coefficients its
signed profile \ref{eq:rational-profile} goes negative, so it fails (ii). Each verdict is a
computation on the printed coefficients of one design and not a theorem about a class of designs.
Whatever their accuracy as Gaussian approximations, a filter whose $R$ fails the test is
the kernel of no exact cascade. The theory's recommendation is therefore the opposite of the
literature's optimization: approximate the Gaussian \emph{from inside the class}, by real-pole
cascades with no zeros, and accept the \Matern{}'s exponential tails and its kink at the origin
for $\gamma = 1$ as the price of exactness in scale.
\end{remark}

% draft: Proposition 14.4, clauses (1) and (2). Clause (3) is NOT stated as a node; see
% rem:rational-converse and the header note.
% CHANGED 2026-09-19 (R185, second external review round, referee C finding 12): FOUR STEPS OF THE
% PROOF OF (2) ARE WRITTEN OUT, none of them new and none of them changing the statement. That
% psi = -log R has no Gaussian coefficient (rational growth against quadratic), so that the pair is
% (0, k_R(x)dx/x) and the substitution u = cx is available; that the exchange of the sum over m with
% the integral is Tonelli on nonnegative summands; that the sum k is the profile of a pair of the
% form eq:sd-profile at all, which is the finiteness of F(1); and that the "exactly when" reads the
% uniqueness of the \Levy{} pair in the direction that forces monotonicity, up to a null set. The
% clause says "with folded \Levy{} profile k_R", so the existence of that profile is a hypothesis
% and is not derived here; the proof now says so.
\begin{proposition}[rational increments and the pyramid they interpolate]
\label{prop:rational-increments}
\uses{prop:thorin-subclass,prop:matern-cascade,thm:main-characterization,lem:selfdecomposable-exponents,lem:quadratic-growth,prop:knot-exactness}
\notes{matern-pyramid}
\begin{enumerate}
\item If
\[
  F(\omega) = \sum_{i=1}^N w_i\log\Bigl(1 + \frac{\omega^2}{\theta_i^2}\Bigr)
\]
with positive integer weights $w_i$ and no Gaussian part, then every increment is rational,
\[
  \hat\mu_{s,t} = \prod_i\Bigl(\frac{1 + s^2\omega^2/\theta_i^2}{1 + t^2\omega^2/\theta_i^2}\Bigr)^{w_i},
\]
of order $2\sum_i w_i$, realized by $\sum_i w_i$ forward and $\sum_i w_i$ backward lag sections,
each pole--zero pair at the common ratio $s/t$.
\item Let $R$ be a symmetric rational function that is the transform of an infinitely divisible
law, with folded \Levy{} profile $k_R$. Then for each $q > 1$ there is exactly one function $F$
that is continuous with $F(0) = 0$ and satisfies $F(q\omega) - F(\omega) = -\log R(\omega)$,
namely
\[
  e^{-F(\omega)} = \prod_{m \ge 1} R\bigl(q^{-m}\omega\bigr),
  \qquad \text{with folded profile} \quad
  k(x) = \sum_{m \ge 1} k_R\bigl(q^m x\bigr);
\]
it is admissible exactly when this $k$ is nonincreasing, in particular whenever $k_R$ is. The
\Matern{} pyramid is the case $R(\omega) = (1+\omega^2)^{-1}$.
\end{enumerate}
\statusT\quad\emph{(The distinction between (1) and (2) is the distinction between the \Matern{}
\emph{family} and the \Matern{} \emph{pyramid}: the family's increments are rational between every
pair of scales, the pyramid's only between scales at ratio $q^n$. So every rational infinitely
divisible transfer generates exactly one pyramid and exactly one candidate fully covariant
family, the interpolation of the pyramid, admissible under a monotonicity condition on its
profile. The index in the profile sum runs from $m = 1$; see the header note.)}
\end{proposition}

\begin{proof}
(1) The increment transform telescopes by \ref{prop:knot-exactness}, and each factor
$\bigl((1+s^2\omega^2/\theta^2)/(1+t^2\omega^2/\theta^2)\bigr)$ is realized as in
\ref{prop:matern-cascade} after the gauge change $t \mapsto t/\theta$; admissibility is
\ref{prop:thorin-subclass} with $U = \sum_i 2w_i\delta_{\theta_i}$.

(2) \emph{Uniqueness.} If $F$ and $\tilde F$ both solve the equation and are continuous with
value $0$ at the origin, then $D := F - \tilde F$ satisfies $D(q\omega) = D(\omega)$, so
$D(\omega) = D(q^{-m}\omega) \to D(0) = 0$ as $m \to \infty$; hence $D \equiv 0$.
\emph{Existence.} Write $\psi := -\log R$, which is a symmetric \Levy{} exponent. Near the origin
$\psi(\omega) = O(\omega^2)$: $R$ is rational and even with $R(0) = 1$, so it is an analytic
function of $\omega^2$ near the origin with value $1$ there, and so is its logarithm, with
value $0$. Then
$\sum_{m \ge 1}\psi(q^{-m}\omega)$ converges, the terms being $O(q^{-2m}\omega^2)$, and its
value $F$ satisfies $F(q\omega) - F(\omega) = \psi(\omega)$ by telescoping. \emph{The pair of
$\psi$.} It has no Gaussian coefficient: $R$ is rational, so $\psi = -\log R$ grows like a
multiple of $\log|\omega|$, while a positive Gaussian coefficient forces quadratic growth. Its
\Levy{} measure is $k_R(x)\,dx/x$ with $k_R$ the folded profile named in the hypothesis, which is
what the substitution below reads; the hypothesis asks for that profile, and
\ref{prop:rational-criterion} exhibits it for every rational $R$ meeting its condition (i) --- it
is the signed profile \ref{eq:rational-profile}, a finite combination of exponentials, and
conditions (ii) and (iii) there are what make it a profile. \emph{The profile.}
Substituting $u = cx$ in the representation shows that $\psi(c\,\cdot)$ has folded profile
$k_R(u/c)$; taking $c = q^{-m}$ and summing over $m \ge 1$ gives
$k(x) = \sum_{m\ge1}k_R(q^mx)$, the exchange of the sum with the integral being Tonelli, the
summands being nonnegative. That $F$ is of the form \ref{eq:sd-profile} at the pair $(0,k)$ ---
that is, that $k$ satisfies the two integrability conditions of \ref{lem:profile-integrability}
--- follows from the finiteness of $F(1)$, which the convergence of the series gives.
\emph{Admissibility} is \ref{thm:main-characterization} through clause (3) of
\ref{lem:selfdecomposable-exponents}, the representation \ref{eq:sd-profile} --- and not that
lemma's differentiability clause, which is a different statement: $F$ is of that form with $k$
nonincreasing exactly when the family is admissible, and $k$ inherits monotonicity from $k_R$
when $k_R$ is nonincreasing, each summand being
so. In the other direction, that admissibility \emph{forces} $k$ to be nonincreasing is the
uniqueness of the \Levy{} pair (\ref{prop:fourier-toolbox}(3)), and it fixes $k$ only up to a null
set, so the monotonicity is monotonicity of a representative.
\end{proof}

% CHANGED (2026-09-19, external review of module B, R170): the remark recorded the converse of
% prop:rational-increments(1) as an expectation ("monotonicity of k should exclude the
% oscillating terms"). The referee's example below refutes it. The remark now states the
% counterexample; the label is kept. Safe direction: an expectation is withdrawn.
% CHANGED 2026-09-19 (R184, second external review round, referee C finding 3): A SECOND EXAMPLE,
% k = 4e^{-x} - 2e^{-2x}, and the case split of rem:recursive-gaussians restated. The first example
% here shows that a complex pole pair survives monotonicity; the referee's shows that the real-pole
% case is not the safe one either -- all poles and zeros real, admissible, outside the Thorin
% subclass (k''(0+) = -4 < 0) and outside the families realized by lag sections, because the stage's
% pole at range s/2 has no smaller zero to pair with. Every printed detail was recomputed here:
% positivity and monotonicity, the exponent by prop:matern-exponent(1) at the rates 1 and 2 with
% weights gamma = 2, theta = 1 and gamma = 1, theta = 1/2, the second derivative, and the pole and
% zero ranges of the stage. So "real poles" implies neither Thorin nor lag sections, and
% rem:recursive-gaussians' first case is restated as real poles AND NO ZEROS. Safe direction: a
% remark's case split is narrowed and a counterexample added.
\begin{remark}[the converse fails: a rational family with a complex pole pair]
\label{rem:rational-converse}
The converse of \ref{prop:rational-increments}(1) is false: an admissible family all of whose
increments are rational, of bounded order, need not have an exponent of the form stated there,
and need not lie in the Thorin subclass. Take $a = 0$ and
\[
  k(x) = 2e^{-x}\,(3 + 2\cos x), \qquad x > 0 .
\]
It is positive, and $k'(x) = -2e^{-x}\,(3 + 2\cos x + 2\sin x) < 0$ because
$2\cos x + 2\sin x \ge -2\sqrt2 > -3$, so $k$ is nonincreasing; both integrability conditions
are immediate. The identity of \ref{prop:matern-exponent}(1) holds for a complex rate of positive
real part --- \ref{eq:complex-frullani}, proved at \ref{prop:rational-criterion}, where it is
first needed --- and the two rates
$1 \pm i$ of the oscillating term combine to a real exponent:
\[
  F(\omega) = 3\log(1 + \omega^2) + \log\Bigl(1 + \frac{\omega^4}{4}\Bigr),
  \qquad
  e^{-F(\omega)} = \frac{1}{(1+\omega^2)^3\,(1 + \omega^4/4)} .
\]
So every increment is rational, of order at most ten. But
$k^{(4)}(x) = 2e^{-x}\,(3 - 8\cos x)$ is negative near the origin, so $k$ is not completely
monotone and the family is not a Thorin member (\ref{prop:thorin-subclass}). Monotonicity of the
profile therefore does not exclude oscillating terms, a real pole of sufficient weight being
enough to carry a complex pair. Such a family has rational stages and no realization by the
lag sections of \ref{prop:matern-cascade}, each of which is positive by itself: the factor
of the complex pair, $1/(1 + \omega^4/4)$, is the transform of a kernel proportional to
$e^{-|x|}(\cos x + \sin|x|)$, which changes sign, so the positivity of the member comes from
the product only.

Complex poles are not what fails, and a second example shows it. Take $a = 0$ and
\[
  k(x) = 4e^{-x} - 2e^{-2x}, \qquad x > 0 .
\]
It is positive, since $e^{-x} \le 1 < 2$; it is nonincreasing, since
$k'(x) = -4e^{-x}(1 - e^{-x}) \le 0$; and both integrability conditions are immediate, $k$ being
bounded and exponentially decaying. The identity of \ref{prop:matern-exponent}(1) at the two rates
gives
\[
  F(\omega) = 2\log(1 + \omega^2) - \log\Bigl(1 + \frac{\omega^2}{4}\Bigr),
  \qquad
  e^{-F(\omega)} = \frac{1 + \omega^2/4}{(1+\omega^2)^2},
\]
so every increment is rational, of order at most six, and every pole and every zero of the
kernel's transform is on the imaginary $\omega$-axis: the transfer has no oscillating factor at
all. Yet $k''(x) = 4e^{-x} - 8e^{-2x}$ is negative at the origin, so $k$ is not completely monotone
and the family is not a Thorin member, hence not in the class of \ref{prop:rational-closure}.
Nor is it realized by lag sections: the stage from $s$ to $t$ has pole ranges $t, t, s/2$ and zero
ranges $t/2, s, s$, and the pole at range $s/2$ has no zero of smaller range to pair with, so any
factorization into first-order sections has a section with a negative coefficient. So ``real
poles'' implies neither membership in the Thorin subclass nor a realization by lag sections, and
the real-pole case of \ref{rem:recursive-gaussians} is the case of real poles \emph{and no zeros}.
\end{remark}

% draft: Proposition 14.5
% twin: hcs prop:thorin-subclass (its implementation clause)
% CHANGED (2026-09-11, draft sync, R82): the rounding clause said "rounding the weights to
% positive integers ... makes every increment rational, realized by sum_i w_i ... sections".
% The weights are those of U, and eq:thorin carries the factor 1/2: U_N = sum_i w_i delta_{theta_i}
% gives F_N = sum_i (w_i/2) log(1 + omega^2/theta_i^2), so an odd w_i leaves a half-integer
% power in the increment, which is not rational. Narrowed to even weights w_i = 2 n_i, realized
% by sum_i n_i sections -- the draft's own form. The proof carried the other normalization
% ("summand w_i log(1+omega^2/theta_i^2) with profile 2 w_i e^{-theta_i x}") and is corrected to
% the one the statement and the parallel generator use; the generator, (1/t) sum_i w_i (Lap - I),
% is prop:corner-generators(4) at U_N and was right. rem:costs is recounted the same way. Found by
% the draft-sync editor of group D (diff row D-4).
% CHANGED (2026-09-19, article mirror of module B, R174): the title and one sentence, from the
% register pass's item 8. The title was "the Thorin subclass as a machine", which the paper replaced
% on 2026-09-15 by "finite Thorin approximations and their realization", the two things the node
% states; the blueprint follows, so that the shared statement and its copy carry one name. The
% closing sentence's circuit metaphor ("The kernels compose in series and the generators add in
% parallel ... the parallel sum ... of a Laplacian and bounded first-order blocks") is replaced by
% what it means: the Thorin sum lives in the exponent, so the kernels compose as a product of
% factors and the generators add, and the summands are a Laplacian and finitely many bounded
% operators. No clause changes. The pass's other findings at this node are moot: "quadrature" is
% defined in the statement itself; the convergence clause is stated at the exponents and needs no
% hypothesis on U_N beyond that; and the convergence the rounding clause does NOT claim is now
% prop:rational-closure (R169, 2026-09-19).
\begin{proposition}[finite Thorin approximations and their realization]\label{prop:thorin-machine}
\uses{prop:thorin-subclass,prop:rational-increments,prop:choquet-cone,prop:levy-continuity,prop:corner-generators,lem:admissible-cone,prop:scale-evolution}
\notes{exact-cascade-implementation}
Let $F$ be admissible with completely monotone profile and Thorin measure $U$ as in
\ref{eq:thorin}. Every measure $U_N = \sum_{i \le N} w_i\delta_{\theta_i}$ with positive weights gives,
with the same Gaussian coefficient, an admissible $F_N$. A sequence of such measures is a
\emph{quadrature} of $U$ when $F_N(\omega) \to F(\omega)$ for every $\omega$, and the kernels of
a quadrature converge weakly to those of $F$ at every scale. \emph{Every such $U$ has a
quadrature}, and it has one in the variant that carries the mass above the last node as a Gaussian
coefficient: there are cutoffs $\Theta_N \uparrow \infty$ and finite measures
$U_N = \sum_i w_i\delta_{\theta_i}$ with positive weights, carried by $(0,\Theta_N]$, such that
both the exponents $F_N$ at Gaussian coefficient $a$ and the exponents $\tilde F_N$ at Gaussian
coefficient
\[
  a_N \ :=\ a + \tfrac12\int_{(\Theta_N,\infty)}\theta^{-2}\,U(d\theta) \ <\ \infty
\]
converge to $F(\omega)$ at every $\omega$; $\tilde F_N$ is admissible as well. Rounding the weights to even
integers $w_i = 2n_i$ keeps $F_N$ admissible and makes the increments of its jump part rational,
realized by $\sum_i n_i$ forward--backward lag sections per knot, in series; a Gaussian part is
a separate factor $e^{-a(t^2-s^2)\omega^2}$ of every increment, and is not rational. The Thorin sum
lives in the exponent, so the stages compose as a product of factors and the generators add: the
generator of \ref{prop:scale-evolution} is the sum
\[
  \mathcal{A}_t = 2at\,\partial_x^2
   + \frac1t\sum_i w_i\bigl(\Lap_{t/\theta_i} - I\bigr)
\]
of a Laplacian and finitely many bounded operators, one per atom.
\statusT\quad\emph{(The constants in the parallel generator are those of
\ref{prop:corner-generators}(4): the Gaussian term carries the factor $t$ and the sum the factor
$1/t$. The draft carried different constants until the draft sync of 2026-09-11; see the header
note. The weights are those of $U_N$, so with the factor $\tfrac12$ of \ref{eq:thorin} the
exponent weights are $w_i/2$, which is why the rounding is to even integers.
\textbf{Existence} (2026-09-19, R199, the author's decision D3). Until today the node
\emph{defined} a quadrature by its convergence and showed that one converges weakly, and never
said that any exists; five sentences of the article, and the numerical run, read it as if it did
(referee C finding 6). The existence clause and its three-line proof are now here. The variant
beside it --- the mass above the last node carried as a Gaussian coefficient rather than as an
atom --- is what the article's numerical run actually computes, and the node did not cover it,
"with the same Gaussian coefficient" being part of the first sentence. \textbf{Not claimed}: any
rate, any optimality of the node placement, and nothing about the \emph{rounded} sequence, which
does not converge in general (\ref{prop:rational-closure}).)}
\end{proposition}

\begin{proof}
Admissibility of $F_N$ is closure of the cone under positive sums
(\ref{lem:admissible-cone}, \ref{prop:choquet-cone}), each summand
$\tfrac12 w_i\log(1+\omega^2/\theta_i^2)$ being admissible with profile $w_ie^{-\theta_ix}$; the
same closure, with the Gaussian ray, gives admissibility of $\tilde F_N$. Weak
convergence of a quadrature is the convergence of the exponents that defines it, together with
\ref{prop:levy-continuity}. Even weights $w_i = 2n_i$ make the exponent
$\sum_i n_i\log(1+\omega^2/\theta_i^2)$, so every increment is rational and has the
realization by \ref{prop:rational-increments}(1). The generator is
\ref{prop:corner-generators}(4) with $U = \sum_i w_i\delta_{\theta_i}$, plus the Gaussian term
of \ref{prop:corner-generators}(1).

\emph{A quadrature exists.} Three facts about \ref{eq:thorin}'s condition are used, and all three
read the comparison $\log(1 + \lambda x) \le \lambda\log(1 + x)$, valid for $\lambda \ge 1$ and
$x \ge 0$ because the two sides agree at $x = 0$ and the derivative of the left is
$\lambda/(1+\lambda x) \le \lambda/(1+x)$: that $U$ is finite on $(0,\Theta]$ for every $\Theta$,
since $\log(1+\theta^{-2}) \ge \log(1+\Theta^{-2}) > 0$ there; that
$\int_{(\Theta,\infty)}\theta^{-2}U(d\theta) < \infty$, since
$\log(1+\theta^{-2}) \ge \tfrac12\theta^{-2}$ for $\theta \ge 1$; and that
$\theta \mapsto \log(1+\omega^2/\theta^2)$ is $U$-integrable at each fixed $\omega$, being at most
$\max(1,\omega^2)\log(1+\theta^{-2})$.

Take $\Theta_N := N$. Partition $[1/N, N]$ into finitely many cells of width $\delta_N$, and let
$U_N$ put the mass $U(\text{cell})$ at a point of each nonempty cell; the weights are positive and
finitely many. On $[1/N,N]$,
\[
  \Bigl|\frac{\partial}{\partial\theta}\log\Bigl(1 + \frac{\omega^2}{\theta^2}\Bigr)\Bigr|
   = \frac{2\omega^2}{\theta(\theta^2+\omega^2)} \ \le\ \frac 2\theta \ \le\ 2N ,
\]
a bound free of $\omega$, so moving each cell's mass to a point of that cell changes the Thorin
integral by at most $2N\delta_N\,U([1/N,N])$, and $\delta_N$ may be chosen to make that at most
$1/N$. What is left is the mass of $U$ off $[1/N,N]$, whose contribution
$\tfrac12\int_{(0,1/N)\cup(N,\infty)}\log(1+\omega^2/\theta^2)\,U(d\theta)$ tends to $0$ at each
fixed $\omega$ by dominated convergence, the integrand being dominated by the $U$-integrable
$\max(1,\omega^2)\log(1+\theta^{-2})$ and the domain shrinking to the empty set. So
$F_N(\omega) \to F(\omega)$ for every $\omega$, which is a quadrature. For the variant, the
difference $\tilde F_N - F_N$ is
$\bigl(\tfrac12\int_{(N,\infty)}\theta^{-2}U(d\theta)\bigr)\omega^2$, and the integral is the tail
of a finite one, so it tends to $0$; hence $\tilde F_N(\omega) \to F(\omega)$ as well.
\end{proof}

The spectra of the working families are those of the causal theory under the bijection
$\theta \mapsto \sqrt{2\theta}$ of \ref{prop:thorin-subclass}(3): the \Matern{} one atom, the
symmetric stable a power law, and, under the hypothesis of \ref{rem:student-ggc}, the Student-t
the Bessel spectral density. For a power law,
logarithmically spaced nodes are the natural quadrature, and the Thorin structure guarantees
that a positive fit exists, positivity of the weights being membership in the axiom class.

% NEW NODE (2026-09-19, the author's decision at the review of module B; fidelity row R169).
% The numerical run of the article (paper-b, ex:numerical-run, 2026-09-18) found that the
% families with rational stages do not approximate the Poisson member: nodes of weight 2 at the
% midpoints of the cells of its Thorin measure (2/pi) d theta converge to the member with
% transform 1/cosh(omega), at L1 distance 0.507, and no placement of such nodes came closer than
% 0.200. prop:thorin-machine never claimed otherwise -- its convergence clause is for the
% quadrature, its rounding clause asserts admissibility and rationality only -- but the prose
% around it, here (rem:cannot-march, rem:costs) and in the article, read as if the rounded
% families converged. This node states what the rounding can reach. It is [A] on the new ledger
% entry A24 (Bondesson's closure theorem for generalized gamma convolutions, with Feller's
% continuity theorem for Laplace transforms), proved in prose; no Lean is scheduled, as for the
% rest of the chapter.
% CHANGED 2026-09-19 (R186, second external review round, referee C finding 2): THE PREAMBLE'S
% IDENTIFICATION IS PROVED IN BOTH DIRECTIONS. It read "These are the families realized by lag
% sections", of which only the inclusion R -> realized was justified (by
% prop:rational-increments(1)), while the proposition's title, rem:costs and rem:cannot-march all
% use the converse. The sense of "lag section" is now fixed at prop:matern-cascade's -- pole and
% zero at the two scales of the stage -- and the converse is the one line the referee asked for: at
% s = 0 the zero is gone, so the kernel from the origin has transform prod (1 + t^2 omega^2 /
% theta_j^2)^{-n_j} and the exponent is in R, which is dilation-invariant. No clause changes; a
% claim already made is now proved on the page.
\begin{proposition}[the closure of the families realized by lag sections]\label{prop:rational-closure}
\uses{prop:thorin-machine,prop:thorin-subclass,prop:rational-increments,prop:levy-continuity,prop:stable-family,prop:matern-exponent,prop:laplace-uniqueness-locally-finite,prop:moments-tails,prop:student-t}
\notes{exact-cascade-implementation}
Let $\mathcal{R}$ be the set of exponents
$F(\omega) = \sum_{i \le N} n_i\log(1 + \omega^2/\theta_i^2)$ with $N \ge 0$, positive integers
$n_i$ and $\theta_i > 0$. These are exactly the families realized by lag sections in the sense of
\ref{prop:matern-cascade}, that is by cascades of first-order sections whose pole and zero sit at
the two scales of the stage: each member of $\mathcal{R}$ is admissible
with no Gaussian part and has every increment a cascade of such sections, by
\ref{prop:rational-increments}(1); and conversely, an admissible family one of whose kernels from
the origin is such a cascade has its exponent in $\mathcal{R}$.
\begin{enumerate}
\item \emph{(Limits.)} Let $F_n \in \mathcal{R}$, let $F$ be admissible, and suppose that the
kernels of $F_n$ at canonical scale $1$ converge weakly to that of $F$. Then $F$ has a symmetric
Thorin representation \ref{eq:thorin} whose Thorin measure is a sum of atoms of even integer
mass, $U = \sum_i 2n_i\delta_{\theta_i}$ with positive integers $n_i$ and with finitely many
$\theta_i$ in every compact subset of $(0,\infty)$. Its Gaussian coefficient may be positive.
\item \emph{(Every such member is a limit.)} Conversely, every admissible $F$ with a symmetric
Thorin representation of that form, with any Gaussian coefficient $a \ge 0$, is the limit in
that sense of a sequence in $\mathcal{R}$.
\item \emph{(What is not a limit.)} Consequently an admissible $F$ is not such a limit if its
profile is not completely monotone, or if its Thorin measure has a part without atoms or an
atom whose mass is not an even integer. In particular no symmetric stable member with
$0 < \alpha < 2$ is such a limit, and no \Matern{} member of non-integer index.
\item \emph{(The class of limits is closed.)} If every $F_n$ is such a limit, $F$ is admissible,
and the kernels of $F_n$ at canonical scale $1$ converge weakly to that of $F$, then $F$ is such
a limit as well.
\item \emph{(Every limit has light tails.)} If $F$ is such a limit then its displacement $X_t$
has finite moments of every order, and an exponential moment: for every $t > 0$ there is
$\beta > 0$ with $\mathbb{E}e^{\beta|X_t|} < \infty$. Consequently an admissible family with a
divergent absolute moment of some order is not such a limit. No symmetric stable member with
$0 < \alpha < 2$ is one, and no Student-t member is one either --- the second exclusion with no
hypothesis on that family's Thorin measure.
\end{enumerate}
\statusA\quad\ledger{A24}\quad\ledger{A13}\quad\emph{(\textbf{Assignment.} \ledger{A24} carries
two cited facts,
both spent by clause (1) and again by clause (4): the continuity theorem for Laplace transforms,
that pointwise
convergence of the transforms of probability laws on $[0,\infty)$ to a limit that tends to $1$
at the origin is weak convergence to a probability law with that transform; and Bondesson's
closure theorem, that a weak limit of generalized gamma convolutions is one, with the Thorin
measures converging vaguely on $(0,\infty)$. \ledger{A13} carries the moment criterion, spent by
clause (5) twice: at the weight $(|x| \vee 1)^n$ through \ref{prop:moments-tails}(1) for the
moments, and at the submultiplicative weight $e^{\beta|x|}$ for the exponential moment, which is
the entry's own form and which no other node of this article reads. Clauses (2) and (3) spend no
cited fact. What is
NOT carried, and not claimed: a rate, or a distance between a member and the closure, which the
article measures in a numerical example for the Poisson and one Student-t member; the value of
$\beta$ in clause (5), which the proof exhibits as anything below $\theta_{\min}/t$ but which the
statement leaves existential; and any converse to clause (5) --- light tails do not make a member
a limit, the \Matern{} member of non-integer index having every moment and an exponential moment
and being excluded by clause (3).
\textbf{Clauses (4) and (5) were added on 2026-09-19} (R200, the author's decision D3). Clause (4)
was used in the article's numerical run and was a clause of nothing (referee C finding 13). Clause
(5) is the external review's observation (referee C finding 7), and what it buys is that the
Student-t exclusion no longer passes through \ref{rem:student-ggc}: until today every statement
that the Student-t family is outside this closure went through that remark's hypothesis, since it
was read off the family's Thorin measure. It is now read off its moments instead, which
\ref{prop:student-t}(4) gives from the explicit density.)}
\end{proposition}

\begin{proof}
\emph{The identification of $\mathcal{R}$.} One direction is
\ref{prop:rational-increments}(1) at $U = \sum_i 2n_i\delta_{\theta_i}$. For the other, let the
kernel of an admissible family at some scale $t > 0$ be a cascade of lag sections in that sense.
A section of the stage from $s$ to $t$ has transfer
$(1 + s^2\omega^2/\theta^2)/(1 + t^2\omega^2/\theta^2)$, so at $s = 0$ its zero is gone and the
cascade's transfer is $\prod_j(1 + t^2\omega^2/\theta_j^2)^{-n_j}$ with positive integer
multiplicities. Taking logarithms, the exponent at that scale is
$\sum_j n_j\log(1 + t^2\omega^2/\theta_j^2)$, which is a member of $\mathcal{R}$ read at the
ranges $\theta_j/t$; and $\mathcal{R}$ is invariant under $\omega \mapsto \lambda\omega$, so the
family's exponent lies in it whatever the gauge.

(1) \emph{The transforms.} Weak convergence of the kernels gives convergence of their cosine
transforms, $e^{-F_n(\omega)} \to e^{-F(\omega)}$ for every $\omega$, and the limit is positive,
so $F_n(\omega) \to F(\omega)$.

\emph{The causal side.} Write $F_n = \sum_i n_i^{(n)}\log(1 + \omega^2/\theta_{i,n}^2)$ and put
$G_n(\sigma) := F_n(\sqrt{2\sigma})$ for $\sigma \ge 0$. Then
$G_n(\sigma) = \int_{(0,\infty)}\log(1 + \sigma/t)\,V_n(dt)$ with
$V_n := \sum_i n_i^{(n)}\delta_{\theta_{i,n}^2/2}$, and $e^{-G_n}$ is the Laplace transform of the
law $T_n$ of a finite sum of independent gamma variables, of shapes $n_i^{(n)}$ and rates
$\theta_{i,n}^2/2$, since a gamma variable of shape $n$ and rate $t$ has the Laplace transform
$(1 + \sigma/t)^{-n}$. In the parametrization of the cited definition $T_n$ is the generalized
gamma convolution with left extremity $0$ and measure $V_n$, which is finite with compact support
in $(0,\infty)$ and so satisfies the integrability conditions asked there. Put
$G(\sigma) := F(\sqrt{2\sigma})$. Then $G_n \to G$ pointwise on $[0,\infty)$, and $G$ is
continuous with $G(0) = 0$, so $e^{-G(\sigma)} \to 1$ as $\sigma \downarrow 0$. By the
continuity theorem for Laplace transforms the laws $T_n$ converge weakly to a probability law
$T$ on $[0,\infty)$ with Laplace transform $e^{-G}$. By the closure theorem $T$ is a generalized
gamma convolution, say with left extremity $b \ge 0$ and measure $V$, so that
$G(\sigma) = b\sigma + \int_{(0,\infty)}\log(1 + \sigma/t)\,V(dt)$, and $V_n \to V$ vaguely on
$(0,\infty)$.

\emph{A vague limit of integer-valued point measures is one.} Each $V_n$ is a finite sum of atoms
of positive integer mass. Let $I$ be an open interval with compact closure in $(0,\infty)$ and
$V(\partial I) = 0$. Squeezing the indicator of $I$ between continuous functions of compact
support whose integrals against $V$ differ by as little as desired shows $V_n(I) \to V(I)$, so
$V(I)$ is a limit of nonnegative integers and is one. The atoms of $V$ are countably many, so
for every $x \in (0,\infty)$ the intervals $(x - \varepsilon, x + \varepsilon)$ with
$V(\{x \pm \varepsilon\}) = 0$ have $\varepsilon$ arbitrarily small, and along them the integers
$V((x-\varepsilon,x+\varepsilon))$ decrease to $V(\{x\})$ and are therefore eventually equal to
it. So every $x$ has a neighbourhood in which $V$ is an atom of integer mass at $x$, or zero,
and nothing else: $V = \sum_i n_i\delta_{t_i}$ with positive integers $n_i$ and with finitely
many $t_i$ in every compact subset of $(0,\infty)$.

\emph{Back to the line.} Reading $G$ at $\sigma = \omega^2/2$,
\[
  F(\omega) = \tfrac12 b\,\omega^2 + \int_{(0,\infty)}\log\Bigl(1 + \frac{\omega^2}{2t}\Bigr)V(dt)
   = a\,\omega^2 + \tfrac12\int_{(0,\infty)}\log\Bigl(1 + \frac{\omega^2}{\theta^2}\Bigr)U(d\theta),
\]
with $a = \tfrac12 b$ and $U$ twice the image of $V$ under $t \mapsto \sqrt{2t}$, which is the
correspondence of \ref{prop:thorin-subclass}(3). The integrability condition of \ref{eq:thorin} is
the finiteness of the Thorin integral at $\omega = 1$, that is $2\bigl(G(\tfrac12) -
\tfrac12b\bigr)$, and $G$ is finite on $[0,\infty)$. So $F$ has the representation (2) of
\ref{prop:thorin-subclass}, and $U = \sum_i 2n_i\delta_{\sqrt{2t_i}}$.

(2) Let $F(\omega) = a\omega^2 + \sum_i n_i\log(1 + \omega^2/\theta_i^2)$ be of that form, the sum
countable and convergent for every $\omega$ by the integrability conditions of \ref{eq:thorin}.
Put $F_N(\omega) := N\log(1 + a\omega^2/N) + \sum_{i \le N} n_i\log(1 + \omega^2/\theta_i^2)$,
omitting the first term when $a = 0$. Then $F_N \in \mathcal{R}$, its first term being the
\Matern{} exponent of index $N$ and range $\sqrt{a/N}$. Since $N\log(1 + x/N) \to x$ for
$x \ge 0$, $F_N(\omega) \to F(\omega)$ for every $\omega$, so the transforms of the kernels
converge pointwise and the kernels converge weakly by \ref{prop:levy-continuity}.

(3) If $F$ is such a limit, clause (1) gives it a representation \ref{eq:thorin} with a Thorin
measure made of atoms of even integer mass. By \ref{prop:thorin-subclass} its profile is then
completely monotone, and the Thorin measure of an admissible $F$ is unique, by the uniqueness
of the data $(a,k)$ with \ref{prop:laplace-uniqueness-locally-finite}. So an $F$ whose profile
is not completely monotone, or whose Thorin measure is of another kind, is not such a limit.
The symmetric stable member with $0 < \alpha < 2$ has a Thorin measure with a density
(\ref{prop:stable-family}), and the \Matern{} member of index $\gamma$ has the Thorin measure
$2\gamma\,\delta_1$, the exponent $\gamma\log(1+\omega^2)$ being \ref{eq:thorin} at that measure
(\ref{prop:matern-exponent}(1)), whose atom is an even integer only for integer $\gamma$.

(4) By clauses (1) and (2) the limits are exactly the admissible $F$ with a symmetric Thorin
representation $U = \sum_i 2n_i\delta_{\theta_i}$ of the stated kind, at any Gaussian coefficient,
so what has to be shown is that this class is closed under the convergence in question. The
argument of clause (1) shows it, read at a sequence from the class rather than from $\mathcal{R}$:
what it uses of the approximating sequence is that $G_n(\sigma) := F_n(\sqrt{2\sigma})$ is
$b_n\sigma + \int_{(0,\infty)}\log(1+\sigma/t)\,V_n(dt)$ with $b_n \ge 0$ and $V_n$ a measure whose
atoms have positive integer mass, satisfying the integrability conditions of the cited definition,
so that $e^{-G_n}$ is the Laplace transform of a generalized gamma convolution. A member of the
class supplies that too: $b_n = 2a_n$ and $V_n$ is half the image of $U_n$ under
$\theta \mapsto \theta^2/2$, whose atoms have mass $n_i$, and the two integrability conditions
follow from the finiteness of $G_n(1)$ by the three regimes in the proof of
\ref{prop:thorin-subclass}, (1) implies (2). Everything after that --- the continuity theorem, the
closure theorem, the integer-valued vague limit and the return to the line --- reads nothing else.

(5) Let $F$ be such a limit, with Gaussian coefficient $a \ge 0$ and Thorin measure
$U = \sum_i 2n_i\delta_{\theta_i}$ as in clause (1), and let $k$ be its profile. If $U = 0$ the
kernels are Gaussian and every moment and every exponential moment is finite. Otherwise the atoms
have a smallest range. Indeed \ref{eq:thorin}'s condition gives
$\log 2 \cdot U((0,1]) \le \int_{(0,1]}\log(1+\theta^{-2})\,U(d\theta) < \infty$, so $U((0,1])$ is
finite, while every atom has mass at least $2$; there are therefore finitely many atoms in
$(0,1]$, and finitely many in every compact subset of $(0,\infty)$ by clause (1), so
$\theta_{\min} := \min_i\theta_i$ exists and is positive. By \ref{prop:thorin-subclass}, (1) if
and only if (2), the profile is $k(x) = \sum_i 2n_ie^{-\theta_ix}$, finite at $x = 1$; hence for
$x \ge 1$
\[
  k(x) = \sum_i 2n_ie^{-\theta_i}e^{-\theta_i(x-1)}
   \ \le\ e^{-\theta_{\min}(x-1)}\sum_i 2n_ie^{-\theta_i}
   \ =\ k(1)\,e^{\theta_{\min}}\,e^{-\theta_{\min}x} .
\]
Every moment is then finite by \ref{prop:moments-tails}(1), since
$\int_1^\infty x^{n-1}k(x)\,dx$ is dominated by a multiple of
$\int_1^\infty x^{n-1}e^{-\theta_{\min}x}dx$. For the exponential moment, fix $t > 0$ and
$\beta < \theta_{\min}/t$. The \Levy{} measure of $X_t$ is the $t$-dilate of the folded
$k(x)x^{-1}dx$, so, as in the proof of \ref{prop:moments-tails},
$\int_{|x|>1}e^{\beta|x|}\,\nu_2(dx) = \int_{1/t}^\infty e^{\beta t x}\,k(x)\,x^{-1}dx$, which the
bound above makes finite: near infinity the integrand is dominated by a multiple of
$e^{(\beta t - \theta_{\min})x}$, and on the compact stretch $[1/t, 1]$ it is bounded, $k$ being
finite and nonincreasing there. The weight $x \mapsto e^{\beta|x|}$ is submultiplicative, so the
moment criterion applies to it and gives $\mathbb{E}e^{\beta|X_t|} < \infty$.

The consequence is the contrapositive. A symmetric stable member with $0 < \alpha < 2$ has
$\mathbb{E}|X_t|^n = \infty$ for $n \ge \alpha$ (\ref{prop:stable-family}), and the Student-t
member of shape $a$ has $\mathbb{E}|X_t|^n = \infty$ for $n \ge 2a$ (\ref{prop:student-t}(4)); so
neither is a limit of families realized by lag sections. The second exclusion reads the moments of
the explicit density of \ref{prop:student-t}(1) and no property of that family's Thorin measure,
so it is free of the hypothesis of \ref{rem:student-ggc}.
\end{proof}

% NEW 2026-09-19 (R201; the author's decision D4, its second half): the attribution of the limit
% class. The external review found that the class of prop:rational-closure is the symmetric Polya
% frequency functions of order infinity, and that Bondesson states the half-line form of the same
% identification on his p. 36, two pages after the closure theorem ledger A24 cites. Both pages
% were read from the page images by the librarian on 2026-09-19
% (notes/reviews/cone/SOURCES-round2-2026-09-19.md, Parts 1.2 and 1.3), and ledger A21 already
% holds Karlin's Thm. 3.2(a), p. 345, with the class E2* of p. 336. The remark is an attribution
% with the two-line computation that makes it one, and it is SCOPED: the identification is of the
% kernels from the origin and not of the families, a stage between two positive scales not being
% of Schoenberg's form at all; and non-creation of extrema is named as untreated here.
\begin{remark}[the limit class is the symmetric P\'olya frequency functions]
\label{rem:polya-limits}
\uses{prop:rational-closure,prop:thorin-subclass}
The class of limits identified in \ref{prop:rational-closure} is classical under another name. A
probability density $f$ on the line is a \emph{P\'olya frequency function of order $\infty$} when
for every $n$ and all $x_1 < \dots < x_n$, $y_1 < \dots < y_n$ the matrix $(f(x_i - y_j))_{i,j}$
has nonnegative determinant. Schoenberg's representation theorem, as printed in
\sscite{karlin1968total}, Thm.~3.2(a), p.~345, with the class $\mathcal{E}_2^{*}$ of p.~336, says
that $f$ is one exactly when the reciprocal of its two-sided Laplace transform $\varphi$ is the
entire function
\begin{equation}\label{eq:schoenberg-e2}
  \frac{1}{\varphi(s)} = e^{-\gamma s^2 + \delta s}\prod_i (1 + a_is)\,e^{-a_is},
  \qquad \gamma \ge 0, \quad \delta,\ a_i \ \text{real}, \quad \sum_i a_i^2 < \infty,
\end{equation}
with the nondegeneracy $\gamma + \sum_i a_i^2 > 0$; those are the pages ledger A21 anchors, and
Karlin attributes the theorem to Schoenberg (1951) in the notes to his chapter, p.~391. The primary source is \sscite{schoenberg1951polya}, Theorem~I, pp.~333--334, read at the page: it
states the representation for the class of his (3), p.~331, without the normalization to a
density that Karlin's statement makes explicit, and the same paper proves the
variation-diminishing property of these functions.

Read at a member of the class, the two agree. Let $F(\omega) = a\omega^2 +
\sum_i n_i\log(1 + \omega^2/\theta_i^2)$ be as in \ref{prop:rational-closure}(1). By clause (5)
its kernel from the origin at canonical scale $1$ has an exponential moment, so its two-sided
Laplace transform converges on a strip about the imaginary axis and is there the analytic
continuation of the cosine transform: $\varphi(s) = e^{-F(-is)} = e^{as^2}\prod_i
(1 - s^2/\theta_i^2)^{-n_i}$, whence
\[
  \frac{1}{\varphi(s)}
   = e^{-as^2}\prod_i\Bigl[\Bigl(1 + \frac{s}{\theta_i}\Bigr)
       \Bigl(1 - \frac{s}{\theta_i}\Bigr)\Bigr]^{n_i} .
\]
That is \ref{eq:schoenberg-e2} with $\gamma = a$, $\delta = 0$ and the numbers $\pm1/\theta_i$ for
the $a_i$, each with multiplicity $n_i$: the factors $e^{-a_is}$ cancel in pairs, and
$\sum_i a_i^2 = \sum_i 2n_i\theta_i^{-2}$, which \ref{eq:thorin}'s condition makes finite and
which is what makes the product entire. Conversely, if $f$ is a symmetric probability density
whose reciprocal transform is \ref{eq:schoenberg-e2}, then evenness makes the zero set of
$1/\varphi$ symmetric, so $\delta = 0$ and the $a_i$ pair as $\pm1/\theta_j$ with multiplicities
$m_j$; the paired factors collapse to $\prod_j(1 - s^2/\theta_j^2)^{m_j}$, and reading at
$s = i\omega$ gives the cosine transform $e^{-\gamma\omega^2}\prod_j
(1+\omega^2/\theta_j^2)^{-m_j}$, an exponent of the form \ref{eq:thorin} at
$U = \sum_j 2m_j\delta_{\theta_j}$ --- admissible by \ref{prop:thorin-subclass}, with finitely
many $\theta_j$ below any bound because $\sum_j m_j\theta_j^{-2} < \infty$. \emph{So the kernels
from the origin of the members of the class of \ref{prop:rational-closure} are exactly the
symmetric P\'olya frequency functions of order $\infty$}, the nondegeneracy
$\gamma + \sum_ia_i^2 > 0$ corresponding to $F \not\equiv 0$, which (ND) supplies. Bondesson
states the half-line form of the same identification, for Thorin's class on $[0,\infty)$, two
pages after the closure theorem used above (\sscite{bondesson1992generalized}, p.~36): ``Obviously
the [$PF_\infty$]-class equals the subclass of [Thorin's class] for which the $U$-measure is
discrete with atoms with integral mass. The closure theorem for GGC's guarantees in particular
that also the $PF_\infty$-class is closed with respect to weak limits; note that atoms with mass
$1$ cannot be split in the limit.''

\emph{What this does not say.} The identification is of the kernels \emph{from the origin}. A
stage between two positive scales is not of Schoenberg's form at all: the reciprocal of its
transform is $(1+t^2\omega^2)/(1+s^2\omega^2)$ read at imaginary argument, which has poles, while
\ref{eq:schoenberg-e2} is an entire function. Nothing is therefore claimed here about the
\emph{families} --- neither that their stages are P\'olya frequency functions nor that they
diminish variation. And the question the identification suggests, whether the scale space of such
a member creates no new local extremum, is not treated in this module; it is left for further
work.
% TODO(module C): cite the non-creation theorem here once module C is released.
\end{remark}

% CHANGED (2026-09-19, article mirror of module B, R174): the two remarks below get the two pointers
% the register pass's item 18 asked for, and rem:cannot-march glosses "marched" at its first use
% (item 5). The ill-posedness was asserted with nothing behind it and is now read at the member
% a = 1/2, where the equation is the Laplace equation and the instability is Hadamard's; the
% accuracy O(1/gamma) of the real-pole Gaussian was established nowhere and is now the expansion of
% rem:recursive-gaussians' own transfer, at a fixed frequency, which is all that is claimed.
% CHANGED (2026-09-19, R169): the two remarks below read as if the lag-section quadrature of
% prop:thorin-machine converged to the Student-t and stable members. It does not
% (prop:rational-closure): the quadrature that converges has real weights and no lag-section
% realization. rem:cannot-march now says which quadrature is meant, and rem:costs no longer
% prices the stable and Student-t families by "their Thorin quadratures". Safe direction.
% CHANGED (2026-09-19, external review of module B, R170): the remark said that the Student-t
% scale space "cannot be computed by stepping in scale from the signal, only ... by its kernel".
% That is false: prop:knot-exactness is a cascade for every admissible family, and the member
% a = 1/2 steps by the Poisson semigroup. What cannot be marched is the LOCAL equation, whose
% Cauchy problem is ill-posed. The title keeps its claim, which was always about the equation.
% CHANGED 2026-09-19 (R187, second external review round, referee E finding 3.4): THE REMARK IS
% SPLIT INTO WHAT IS ESTABLISHED AND WHAT IS ANNOUNCED. Module B's section 5.4 says that which
% principle cuts the Student-t family out of the cone is the subject of the next module and that
% nothing is stated about it here, while this remark asserted the elliptic boundary-value problem,
% the ill-posedness of its Cauchy problem, and Hadamard's instability. Of those, the equation itself
% is thm:joint-locality -- chapter 12, module C, and machine-checked at the family
% (SpatialLine.student_scaleSpace_gaspt) -- but the ILL-POSEDNESS is asserted by no statement of
% this article at any member, the a = 1/2 reading being classical and the general case a
% description. The remark now says which half is which; nothing proved is withdrawn (the steppable
% cascade, the absent lag-section realization and the two exclusions all stay, with their nodes).
% Safe direction: a remark claims less.
% CHANGED 2026-09-19 (R202, the author's decision at the second review round): the two references
% into chapter 12 are replaced by prose, module C being unpublished and a reference across the
% module boundary being what a repository split would break; and the Student-t exclusion no longer
% passes through rem:student-ggc, prop:rational-closure(5) having made it unconditional (R200).
% TODO(module C): cite thm:joint-locality and eq:gaspt here once module C is released.
\begin{remark}[the local equation of the Student-t corner cannot be marched]\label{rem:cannot-march}
To \emph{march} an equation in scale is to step its solution forward from the signal by a local
update, as an explicit scheme does with the heat equation. On the line the scale-locally generated
corner is the heat equation, marched routinely. The
jointly local corner, the Student-t family, is of a different kind: its scale space solves an
elliptic boundary-value problem in the half-plane of scale and position rather than a parabolic
evolution in scale. That equation is the subject of a further module, on locality and the
selection of the Gaussian, and no statement of this module asserts it.

\emph{What this chapter establishes about it.} The family can be stepped in scale, as every
admissible family can (\ref{prop:knot-exactness}); its stages are nonlocal operators, those of the
Poisson semigroup for the member $a = \tfrac12$. What it lacks is a realization by lag sections:
the Thorin quadrature of \ref{prop:thorin-machine} with the Bessel spectral density has real
weights and therefore no such realization, and by \ref{prop:rational-closure}(5) no family
realized by lag sections converges to it either, its moments being finite only below the order
$2a$ while every such limit has all of them. So the stages of this corner are computed from its
kernels and not from a recursion in $x$.

\emph{What is announced and not proved here.} That the Cauchy problem of that equation in $t$ is
ill-posed, so that the local equation cannot be marched in scale from the signal, is asserted by no
statement of this article; at the member $a = \tfrac12$ the equation is the Laplace equation, whose
solutions at frequency $\omega$ combine $e^{-t|\omega|}$ with $e^{+t|\omega|}$ --- a perturbation of
the data at high frequency amplified without bound, Hadamard's classical instability --- and the
Poisson stages select the decaying solution; for the other members it is a description and not a
claim. Read with that reservation, the parabolic-versus-elliptic asymmetry between the two local
corners has this as its implementation face.
\end{remark}

\begin{remark}[costs]\label{rem:costs}
Per sample and per knot: the \Matern{} family costs $2\gamma$ first-order states and $4\gamma$
multiply--adds --- two per section, one for the state and one for the output, and $2\gamma$
sections, $\gamma$ forward and $\gamma$ backward --- in a forward and a backward pass at that
knot; a Thorin approximant $U_N = \sum_i 2n_i\delta_{\theta_i}$ costs $2\sum_i n_i$ states and
$4\sum_i n_i$ multiply--adds; the Gaussian by
real-pole recursion costs $2\gamma$ states at an accuracy in the exponent that is $O(1/\gamma)$ at each
fixed frequency (the expansion of $\gamma\log(1 + \sigma^2\omega^2/2\gamma)$ against
$\sigma^2\omega^2/2$, the two transfers of \ref{rem:recursive-gaussians}), and by a direct
finite-impulse-response filter of length $L$ costs $L$; the stable and Student-t families have
no such realization and are not limits of families that have one
(\ref{prop:rational-closure}(5), on their moments alone and with no hypothesis on either family's
Thorin measure), so their stages are computed from their kernels or their
transforms, at the cost of the method chosen. The counts above are totals over the ladder; the
stages can be run one after another, so the live state of an implementation is that of one
stage. This article makes no empirical
claim about natural images.
\end{remark}
